\uuid{7Flt}
\exo7id{7542}
\titre{exo7 7542}
\auteur{mourougane}
\organisation{exo7}
\datecreate{2021-08-10}
\isIndication{false}
\isCorrection{false}
\chapitre{Suite et série de fonctions}
\sousChapitre{Convergence simple, uniforme, normale}
\cpge{mpsi-mp2i.nombres-reels-suites.h, pcsi.nombres-reels-suites.f, ptsi.nombres-reels-suites.f, mpsi-mp2i.polynomes-fractions-rationnelles.a, pcsi.polynomes.a, ptsi.polynomes.a, bcpst1.polynomes-reels.a, mp-mpi.suites-series-fonctions.a, pc.suites-series-fonctions.a, psi.suites-series-fonctions.a, mp-mpi.topologie-espaces-normes.b, pc.espaces-vectoriels-normes.b, psi.espaces-vectoriels-normes.b}
\module{Analyse}
\niveau{L2}
\difficulte{}

\contenu{
\question{
Soit $d$ un entier naturel.
Montrer qu'une suite de polynômes $(P_n(z))_{n\in\Nn}$ de degré au plus $d$ converge localement uniformément sur $\Cc$
si et seulement si les $d+1$ suites de ses coefficients convergent,
si et seulement si, il existe $d+1$ points $c_i$ de $\Cc$ tels que les suites $(P_n(c_i))_{n\in\Nn}$ convergent.
}
}
